CI pipelines rerun whole regression suites on every pull request, however narrow the change. Selecting tests with an LLM at pull-request time can spend more on tokens and latency than it saves. We present \tool, which keeps the LLM off the critical path. An LLM reads each test once offline to produce a short description, and at commit time the diff is scored against the embedded descriptions by cosine similarity. A threshold calibrated by weighted conformal risk control partitions the tests, a telemetry guard forces recently failing tests to run, and only tests near the threshold reach a live LLM. On 420 builds of Apache ShardingSphere with 239 failures, the calibrated semantic layer alone reaches 59.41\% fault recall, the guard alone 92.47\%, and both together 99.58\%. The semantic layer recovers 17 failures that have no failure history, at an estimated \$2.04 per failure on a \$0.10/hour runner \Fix{N19: assumes 5\,s per test}. A gradient-boosted ranker reaches the same recall at higher reduction on this project, but cannot be trained on Apache Log4j2, whose history before evaluation has no failures. We report the net cost after tokens and show that the expected risk our calibration targets is not the fault recall a practitioner reads.
